\begin{table}[h]
\centering
\caption{Urban-Rural Distinctions in MM District Compactness}
\label{tab:urban_rural_comparison}
\begin{tabular}{lccc}
\hline
District Type & N & Mean PP & Std PP \\
\hline
Urban MM & 10 & 0.214 & 0.083 \\
Rural MM & 3 & 0.248 & 0.036 \\
Mixed MM & 6 & 0.182 & 0.050 \\
\hline
\textbf{All MM} & 19 & 0.209 & 0.069 \\
\textbf{All Non-MM} & 50 & 0.236 & 0.072 \\
\hline
\end{tabular}
\vspace{0.2cm}
\begin{minipage}{0.9\textwidth}
\footnotesize
\textbf{Notes:} 
Urban MM districts are concentrated in metro areas (Atlanta, Birmingham, New Orleans, Jackson). 
Rural MM districts span agricultural regions and Black Belt counties. 
Mixed MM districts combine urban and rural populations. 
Urban vs Rural difference: $t=-0.65$, $p=0.528$. 
PP = Polsby-Popper score (higher = more compact).
\end{minipage}
\end{table}